\documentclass[sigplan,screen]{acmart}
\let\Bbbk\relax

\renewcommand\footnotetextcopyrightpermission[1]{}
\pagestyle{plain}

\usepackage{xr}
\usepackage{amsmath,amssymb}

\title{Supplementary Materials for: Merkle Tree-Based Inheritance in Decentralized Module Systems}

\begin{CCSXML}
<ccs2012>
   <concept>
       <concept_id>10011007.10011006.10011008</concept_id>
       <concept_desc>Software and its engineering~Software design techniques</concept_desc>
       <concept_significance>500</concept_significance>
   </concept>
   <concept>
       <concept_id>10002944.10011123.10011131</concept_id>
       <concept_desc>General and reference~Evaluation</concept_desc>
       <concept_significance>300</concept_significance>
   </concept>
</ccs2012>
\end{CCSXML}

\ccsdesc[500]{Software and its engineering~Software design techniques}
\ccsdesc[300]{General and reference~Evaluation}

\keywords{Merkle trees, decentralized systems, module inheritance, software design, evaluation}

\begin{document}
\appendix

\section{Generic Types and Structural Constraints}
\label{appendix:generics}

Though type-agnostic by default, the system accommodates ecosystems with parametric polymorphism and structural typing. This section outlines optional extensions for expressing type-level abstraction and constraints across modules, layered atop the Merkle-based model without modifying its core semantics.

\subsection{Generic Type Aliases}

Modules may define parameterized type aliases of the form:
\[
\texttt{type}~T\langle X_1, \ldots, X_n \rangle = \tau
\]
where $\tau$ may reference type parameters $X_i$. When imported, these aliases are instantiated with concrete types:
\[
T\langle \sigma_1, \ldots, \sigma_n \rangle
\]
Instantiations are resolved into fully derived types within the importing module.

\subsection{Structural Constraints}

Type parameters may be constrained by structural predicates:
\[
\texttt{type}~Map\langle K: \texttt{Hashable}, V \rangle = \ldots
\]
Here, \texttt{Hashable} denotes a required property or interface, transitively resolved through the module graph. Finalization ensures each constraint is satisfied by evidence such as declared instances, structural matches, or subtype edges.

\subsection{Finalization Semantics}

Generic types are instantiated and expanded during finalization. Constraint satisfaction is resolved transitively via module inheritance, yielding deterministic, content-addressed types. The result must be fully resolved and non-parametric in the final Merkle tree.

\subsection{Applications}

These mechanisms enable:
\begin{itemize}
  \item Parameterized types (\texttt{List<T>}, \texttt{Map<K, V>}).
  \item Generic algorithms over abstract constraints.
  \item Interface-based modularity without coupling to implementations.
\end{itemize}

Use of generics is optional. Ecosystems lacking type-level abstraction can ignore this layer without affecting correctness or resolution semantics.

\section{Connection to Existing Module Calculi}
\label{appendix:module-calculi}

This appendix generalizes conventional module calculi by embedding deterministic, content-addressable resolution into the module structure itself. It draws conceptual links to higher-order module systems such as those in ML~\cite{ml_modules}, as well as the mixin calculus~\cite{mixin_calculus}, while extending them for decentralized environments.

Unlike traditional systems, this approach enforces consistent inheritance order via Merkle tree construction and cryptographic tracing, providing explicit override histories and resolution paths even in complex graphs.

\paragraph{Future Work}

We anticipate extending this system with richer abstractions such as recursive types, function modules, and higher-order modules. These additions would support more expressive module ecosystems without compromising determinism or verifiability.

\section{Benchmarks and Evaluation Methodology}
\label{appendix:supplementary-evaluation}

This appendix provides full details on the synthetic benchmark methodology, test scenarios, and comparative metrics used to evaluate the Merkle tree-based module system discussed in the main paper.

\subsection{Synthetic Evaluation}

To validate the model’s feasibility, we constructed synthetic benchmarks over simulated module graphs with controlled topologies:

\begin{itemize}
  \item \textbf{Linear chains:} Deep ancestry resolution.
  \item \textbf{Wide graphs:} Shared base modules for deduplication tests.
  \item \textbf{Diamond graphs:} Shadowing and conflict resolution.
\end{itemize}

Each module contains $N$ symbols and inherits from $k$ parents, with $N \in [10, 1000]$ and $k \in [1, 10]$. Benchmarked operations include:

\begin{itemize}
  \item \textbf{Finalization time:} Merkle root and inheritance trace computation.
  \item \textbf{Symbol resolution:} Lookup and latency analysis.
  \item \textbf{Graph compression:} Deduplication via structural sharing.
\end{itemize}

\paragraph{Disclaimer}
These tests were conducted using a simulated interpreter, not a full runtime. Results are indicative of theoretical scalability, not real-world performance.

\subsection{Results Summary}

\begin{itemize}
  \item Finalization scales linearly with symbol count and ancestry depth.
  \item Symbol lookup remains sublinear in common cases due to caching and reuse.
  \item Structural sharing yields 60--80\% graph compression in shared-ancestor cases.
\end{itemize}

\subsection{Prototype Implementation Plan}

A prototype is planned in \texttt{Rust}, using:

\begin{itemize}
  \item \texttt{libp2p} for decentralized lookup and DHT-based resolution.
  \item \texttt{Wasmer} to evaluate WebAssembly-based modules and execute finalization logic.
\end{itemize}

Tangl is one such prototype implementation of this model and will be used to empirically validate performance in future work.

\subsection{Planned Evaluation Categories}

\begin{itemize}
  \item \textbf{Scalability:} Nested resolution, wide fan-out.
  \item \textbf{Fault tolerance:} Stale caches, resolution under failure.
  \item \textbf{Security:} Malicious ancestry, forked resolution graphs.
  \item \textbf{Tooling usability:} CLI/IDE feedback, trace inspection.
  \item \textbf{Workflow validation:} CI rebuilds, development diffs.
\end{itemize}

\subsection{Comparison to Existing Systems}

Tangl will be benchmarked against \texttt{npm}, \texttt{Cargo}, and \texttt{pip} on tasks including:

\begin{itemize}
  \item Cold vs. warm start resolution.
  \item Incremental rebuilds.
  \item Registry fetch latency.
\end{itemize}

Metrics will include:

\begin{itemize}
  \item DHT hops per module.
  \item Cache hit/miss ratios.
  \item Fallback success under partial connectivity.
\end{itemize}

\section{CI/CD and Version Control Integration}
\label{appendix:ci-vcs}

This appendix illustrates how the system can integrate into real-world workflows, using Tangl as an example implementation.

\subsection{CI/CD Pipelines}

Tangl enables efficient CI/CD by leveraging Merkle roots for incremental validation. Because finalization captures full ancestry, CI systems can:

\begin{itemize}
  \item Cache Merkle roots of finalized subgraphs.
  \item Skip recomputation when inheritance trees are unchanged.
  \item Identify resolution-affecting changes early in the build pipeline.
\end{itemize}

\subsection{Version Control and Tagging}

Ancestry DAGs can be embedded in Git tags or commit annotations. This allows:

\begin{itemize}
  \item Cryptographic validation of ancestry consistency.
  \item Detection of semantic changes from resolution graph diffs.
  \item Structural tagging schemes for symbolic reproducibility.
\end{itemize}

\subsection{Tooling Extensions}

Future versions of Tangl may support:

\begin{itemize}
  \item IDE visualizations of symbol overrides.
  \item Git hooks for ancestry change detection.
  \item CI runners that enforce deterministic finalization traces.
\end{itemize}

These features position Tangl as a prototype that demonstrates the viability of secure, decentralized module resolution in modern software development workflows.

\bibliographystyle{ACM-Reference-Format}
\bibliography{bibliography}

\end{document}
